\section{Background}
\label{sec:background}

\subsection{ILP Formulation Theory}

Integer Linear Programming extends linear programming by requiring some
or all decision variables to take integer values.
The standard reference is Nemhauser and Wolsey~\citeyear{nemhauser1988},
whose Chapter~1 establishes that the feasible region of an ILP is the
convex hull of the feasible integer points, and that branch-and-bound
with LP relaxation bounds provides an exact algorithm with finite
termination.

For a minimisation ILP with objective $\min c^\top x$ subject to
$Ax \leq b$, $x \in \{0,1\}^n$, the branch-and-bound algorithm
maintains an upper bound $U$ (the best feasible integer solution found
so far) and a lower bound $L$ (the LP relaxation at each node).
When $L = U$, the current solution is proven optimal and the algorithm
terminates with a certificate.
The \emph{optimality gap} at any point is $(U - L)/U$; most practical
solvers accept a user-specified gap threshold (e.g., 1\%) rather than
insisting on exact optimality, which can reduce solve time significantly
on larger instances.

\subsection{MTZ vs.\ DFJ Contiguity Constraints}

Enforcing geographic contiguity --- the requirement that each district
forms a connected subgraph --- is the central challenge of redistricting
ILP formulations.
Two constraint families dominate the literature:

\paragraph{Miller--Tucker--Zemlin (MTZ)~\citep{mtz1960}.}
Originally developed for the Travelling Salesman Problem, MTZ constraints
introduce auxiliary flow variables on directed edges.
Each district $d$ designates a root $r_d$; flow from $r_d$ to each
non-root tract enforces a spanning tree over the district's tracts.
MTZ constraints add $O(|E| \cdot k)$ continuous variables and
$O(|E| \cdot k)$ constraints, which is manageable for small instances.
The LP relaxation is relatively weak (the LP bound can be far from the
true integer optimum), but the simplicity of implementation and the
absence of cut generation make MTZ attractive for Phase~1.

\paragraph{Dantzig--Fulkerson--Johnson (DFJ).}
DFJ constraints impose connectivity directly via subtour-elimination
cuts: for every subset $S \subsetneq V_d$ of district $d$'s tracts,
at least one edge must cross the boundary of $S$.
DFJ yields a much tighter LP relaxation than MTZ, which typically reduces
branch-and-bound tree size substantially.
The drawback is that DFJ requires exponentially many constraints
(one per subset), necessitating lazy cut generation at each branch node
(a branch-and-cut implementation).
DFJ is deferred to Phase~2; Phase~1 uses MTZ for implementation simplicity.

\subsection{Both-Endpoint Capacity Bound}

Standard MTZ uses the capacity bound $f_{t,t',d} \leq (n-1) \cdot x_{t,d}$,
which allows flow on edge $(t, t')$ as long as the source $t$ is in
district $d$.
This permits ``spurious'' flow paths: flow can traverse edge $(t, t')$
even if $t'$ is not in district $d$, if $x_{t,d} = 1$.

The \emph{both-endpoint} bound tightens this to
\[
  f_{t,t',d} \leq (n-1) \cdot \min(x_{t,d},\, x_{t',d}),
\]
which forces flow to zero if either endpoint is outside district $d$.
This yields a stronger LP relaxation than standard MTZ at the cost of
replacing the linear bound with a bilinear term (which is linearised
in the ILP using the standard $\min(a,b) \leq a$, $\min(a,b) \leq b$
encoding).
Gurnee and Shmoys~\citeyear{gurnee2021} identify the both-endpoint bound
as the key practical improvement over standard MTZ for redistricting
instances; we adopt it as the default.

\subsection{Prior Work on ILP Redistricting}

Gurnee and Shmoys~\citeyear{gurnee2021} introduce Fairmandering, a
column-generation heuristic for fairness-optimised districting.
Their Section~2 surveys ILP redistricting formulations and identifies
the both-endpoint MTZ bound as the most effective contiguity constraint
for instances with $n \leq 500$ tracts.
They report solve times of 30--120\,s for $k=2$ instances with $n=250$
on a 2019 laptop, consistent with the empirical results in
Section~\ref{sec:results} of this paper.

King, Jacobson, and Sewell~\citeyear{king2015} study contiguity and
hole detection in geographic partitioning and provide the theoretical
foundation for why connected-district constraints are necessary:
without them, ILP solutions may produce district fragments that are
legally invalid and computationally arise naturally at LP relaxation
boundaries.

\subsection{Comparison with Heuristic Approaches}

Table~\ref{tab:method-compare} places \ilp{} in the context of existing
structure-layer algorithms.

\begin{table}[t]
\centering
\caption{%
  Structure-layer algorithm comparison. \ilp{} is the only method providing
  an optimality certificate; tractability is limited to small instances.
}
\label{tab:method-compare}
\renewcommand{\arraystretch}{1.3}
\begin{tabular}{@{}llllll@{}}
\toprule
Property & \ilp{} (B.24) & \metis{} & \sa{} (B.19) & \cvd{} (B.22) & \bfsg{} (B.23) \\
\midrule
Optimality    & Exact cert. & Heuristic & Heuristic & Heuristic & None \\
Contiguity    & MTZ (exact) & Post-proc. & Maintained & Post-proc. & BFS (exact) \\
Tractability  & $n \leq 500$ & Any & Any & Any & Any \\
Runtime       & Exp.\ worst & $O(n \log n)$ & $O(T n)$ & $O(I m)$ & $O(n \log n)$ \\
EC quality    & Gold std.   & High & High & Med & Baseline \\
PP quality    & Depends     & High & Med & High & Med-High \\
Use case      & Cert., small & Prod. & Large & Coastal & Warm-start \\
\bottomrule
\end{tabular}
\end{table}

The key distinction is not runtime but \emph{what the method can assert}.
\metis{}, SA, and CVD can assert ``this plan has EC $= k$''; only \ilp{}
can assert ``\emph{no valid plan} has EC $< k$.''
This distinction is legally material in contested redistricting litigation,
where opposing counsel can always argue that heuristic plans leave
compactness on the table.
